\documentclass[11pt, a4paper]{article}

% ── Packages ──────────────────────────────────────────────────────────
\usepackage[utf8]{inputenc}
\usepackage[T1]{fontenc}
\usepackage{lmodern}
\usepackage[margin=1in]{geometry}
\usepackage{amsmath, amssymb, amsthm}
\usepackage{booktabs}
\usepackage{longtable}
\usepackage{graphicx}
\usepackage{hyperref}
\usepackage{xcolor}
\usepackage{listings}
\usepackage{tcolorbox}
\usepackage{polya}

\hypersetup{
  colorlinks=true,
  linkcolor=polyablue,
  urlcolor=polyablue,
  citecolor=polyablue,
}

% ── Code listing style ────────────────────────────────────────────────
\lstdefinestyle{polya-python}{
  language=Python,
  basicstyle=\ttfamily\small,
  keywordstyle=\color{polyablue}\bfseries,
  commentstyle=\color{polyagray}\itshape,
  stringstyle=\color{polyagreen},
  numbers=left,
  numberstyle=\tiny\color{polyagray},
  numbersep=8pt,
  frame=single,
  framerule=0.4pt,
  rulecolor=\color{polyagray},
  breaklines=true,
  showstringspaces=false,
  tabsize=4,
}
\lstset{style=polya-python}
\title{Feature Importance \& Model Comparison}
\author{uber-polya}
\date{2026-02-26}

\begin{document}

\maketitle
\tableofcontents
\newpage

% ── Phase A: Problem Formulation ─────────────────────────────────────
\section{Problem Formulation}

\begin{polyamodel}

\textbf{Problem Type}: Find \hfill
\textbf{Domain}: Machine Learning


\end{polyamodel}

% ── Universe ──────────────────────────────────────────────────────────
\subsection{Universe}

\begin{itemize}
  \item $Pca: \{explained_variance_ratio, cumulative_variance, n_components_95pct\}$
  \item $Mutual Information: \{feature_scores, feature_ranking, selected_indices\}$
  \item $Rfe: \{feature_ranking, selected_indices\}$
  \item $Feature Overlap: \{mi_rfe_overlap, overlap_count\}$
  \item $Best Combination: \{model_name, feature_method, accuracy, accuracy_std\}$
\end{itemize}

% ── Variables ─────────────────────────────────────────────────────────
\subsection{Variables}

\begin{longtable}{lll}
\toprule
\textbf{Symbol} & \textbf{Meaning} & \textbf{Type / Range} \\
\midrule
\endhead
$x$ & decision variable & $\mathbb{{R}}$ \\
\bottomrule
\end{longtable}

% ── Structure ─────────────────────────────────────────────────────────
\subsection{Structure}

- **Data**: 800 samples, 20 features, binary target (\textasciitilde{}balanced classes)
- **True structure**: Only 5 features are informative, 3 are redundant, 12 are noise
- **Questions**: (1) How many components explain 95\% of variance? (2) Which features matter most? (3) Which model + feature set combination is best?

% ── Mapping ───────────────────────────────────────────────────────────
\subsection{Real-World Mapping}

\begin{longtable}{ll}
\toprule
\textbf{Real-World Concept} & \textbf{Mathematical Object} \\
\midrule
\endhead
problem input & $instance parameters$ \\
\bottomrule
\end{longtable}

% ── Constraints ───────────────────────────────────────────────────────
\subsection{Constraints}

\begin{enumerate}
  \item $PCA explained variance**: How much total variance each principal component captures; cumulative sum determines dimensionality reduction threshold$
  \hfill \textit{()}
  \item $Mutual information**: Information-theoretic measure of dependency between a feature and the target; zero means independence$
  \hfill \textit{()}
  \item $Recursive Feature Elimination (RFE)**: Iteratively removes features with the smallest importance weights from a fitted model$
  \hfill \textit{()}
  \item $Pipeline**: Standardized workflow of preprocessing, feature selection, and model fitting$
  \hfill \textit{()}
  \item $Cross-validation (5-fold)**: Splits data into 5 folds, trains on 4, tests on 1, rotates; provides unbiased accuracy estimate with variance$
  \hfill \textit{()}
  \item $Model comparison**: Evaluating multiple classifiers on identical data splits to identify the best approach for a given dataset$
  \hfill \textit{()}
\end{enumerate}

% ── Objective / Claim ─────────────────────────────────────────────────
\subsection{Claim}


% ── Approach ──────────────────────────────────────────────────────────
\subsection{Approach}

\begin{description}
  \item[Suggested approach:] PCA + Mutual Information + RFE (Random Forest) + 4-Model 5-Fold CV
  \item[Complexity class:] 
  \item[Available tools:] PCA + Mutual Information + RFE
\end{description}
% ── Phase B: Solution ────────────────────────────────────────────────
\section{Solution}

\begin{polyaresult}

\textbf{Answer}: Solution computed


\textbf{Optimal}: No / Unknown \hfill
\textbf{Feasible}: No
\textbf{Algorithm}: PCA + Mutual Information + RFE (Random Forest) + 4-Model 5-Fold CV \quad $$

\textbf{Time}: 5.616892437974457s


\end{polyaresult}

% ── Solution Details ──────────────────────────────────────────────────
\subsection{Solution Details}


Method: 1. PCA (Principal Component Analysis): Fit on all 20 features. Record explained variance per component and find the number of components needed for 95\% cumulative variance. Reveals how much redundancy exists in the feature space.

2. Mutual Information feature selection: Score each feature's mutual information with the target. Rank all 20 features and select the top 8. Filter method -- fast, model-agnostic, detects nonlinear relationships.

3. Recursive Feature Elimination (RFE): Use a Random Forest estimator to recursively eliminate the least important features. Select the top 8. Wrapper method -- slower but accounts for feature interactions through the model.

4. Model comparison (5-fold CV): Compare Logistic Regression, Random Forest, SVM (RBF kernel), and k-NN on three feature sets: all 20 features, MI top 8, and RFE top 8. Each combination evaluated with 5-fold stratified cross-validation accuracy.

5. Best combination: Identify the model + feature set with the highest mean CV accuracy.

% ── Verification ──────────────────────────────────────────────────────
\subsection{Verification}

\begin{longtable}{lcl}
\toprule
\textbf{Check} & \textbf{Status} & \textbf{Detail} \\
\midrule
\endhead
Pca 95Pct Lt 20 & \passmark & Passed \\
Informative Features Selected & \passmark & Passed \\
Rfe Features Selected & \passmark & Passed \\
Selected Beats Random & \passmark & Passed \\
Best Accuracy Gt 80Pct & \passmark & Passed \\
Cv Scores Stable & \passmark & Passed \\
All Passed & \passmark & Passed \\
\bottomrule
\end{longtable}

% ── Phase C: Interpretation ──────────────────────────────────────────
\section{Interpretation}

% ── The Question & The Answer ─────────────────────────────────────────
\subsection{The Question}
- **Data**: 800 samples, 20 features, binary target (\textasciitilde{}balanced classes)
- **True structure**: Only 5 features are informative, 3 are redundant, 12 are noise
- **Questions**: (1) How many components explain 95\% of variance? (2) Which features matter most? (3) Which model + feature set combination is best?.

\subsection{The Answer}

\begin{polyainsight}[title={Bottom Line}]
A feasible solution was found.
\end{polyainsight}

% ── What This Means ───────────────────────────────────────────────────
\subsection{What This Means}

- PCA: Fewer than 20 components needed for 95\% variance (dimensionality is compressible)
- Mutual information: Top 8 features recover at least 3 of the 5 true informative features
- RFE (Random Forest): Top 8 features also recover at least 3 of the 5 true informative features
- Best model accuracy: > 80\% via 5-fold cross-validation
- Feature selection helps: Best selected-feature accuracy >= all-features accuracy - 0.05
- MI/RFE overlap: Substantial agreement between the two selection methods
- All 7 verification checks pass

1. PCA (Principal Component Analysis): Fit on all 20 features. Record explained variance per component and find the number of components needed for 95\% cumulative variance. Reveals how much redundancy exists in the feature space.

2. Mutual Information feature selection: Score each feature's mutual information with the target. Rank all 20 features and select the top 8. Filter method -- fast, model-agnostic, detects nonlinear relationships.

3. Recursive Feature Elimination (RFE): Use a Random Forest estimator to recursively eliminate the least important features. Select the top 8. Wrapper method -- slower but accounts for feature interactions through the model.

4. Model comparison (5-fold CV): Compare Logistic Regression, Random Forest, SVM (RBF kernel), and k-NN on three feature sets: all 20 features, MI top 8, and RFE top 8. Each combination evaluated with 5-fold stratified cross-validation accuracy.

5. Best combination: Identify the model + feature set with the highest mean CV accuracy.

% ── Sensitivity Analysis ──────────────────────────────────────────────

% ── Figures ───────────────────────────────────────────────────────────

% ── Recommendations ───────────────────────────────────────────────────
\subsection{Recommendations}

\begin{enumerate}
  \item Review the model assumptions to ensure real-world fidelity.
\end{enumerate}

% ── Limitations ───────────────────────────────────────────────────────
\subsection{Limitations}

\begin{polyawarnbox}
\begin{itemize}
  \item Model assumes deterministic parameters; real-world variability not captured.
  \item Solution quality depends on the accuracy of input data.
\end{itemize}
\end{polyawarnbox}

% ── Appendix ─────────────────────────────────────────────────────────

\end{document}